% Vector search—the task of finding the  k-nearest neighbors of high-dimensional vectors—is fundamental to machine learning applications such as recommendation systems and information retrieval. However, existing approximate nearest neighbor (ANN) methods struggle under dynamic and skewed workloads common in real-world scenarios, where data and query distributions evolve over time. We introduce Quake, an adaptive indexing system for online vector search that maintains low latency and high recall in such challenging environments. Quake employs an adaptive hierarchical partitioning scheme that adjusts to changing data and access patterns, optimizing performance without the heavy overhead associated with updating graph structures. It incorporates a novel cost model for query latency that predicts performance based on partition sizes and access frequencies, guiding efficient index maintenance. Finally, Quake features adaptive partition scanning to adjust query parameters on-the-fly for meeting recall targets as the index evolves. Further, Quake utilizes NUMA-aware intra-query parallelism to maximize memory bandwidth on multi-core machines. To evaluate Quake, we prepared a Wikipedia vector search workload derived from publicly available data, and developed a workload generator for creating vector search workloads with configurable query and update patterns. On the Wikipedia workload our results show that Quake reduces query latency by $1.5-8\times$ and update latency by $6-83\times$ compared to graph indexes, namely, DiskANN and HNSW. Quake also reduces query latency by $5-22\times$ compared to partitioned indexes, namely, SCANN and Faiss-IVF across all evaluated workloads. 

Vector search, the task of finding the k-nearest neighbors of a query vector against a database of high-dimensional vectors, underpins many machine learning applications, including retrieval-augmented
generation, recommendation systems, and information retrieval. However, existing approximate nearest neighbor (ANN) methods perform poorly under dynamic and skewed workloads where data distributions evolve. We introduce Quake, an adaptive indexing system that maintains low latency and high recall in such environments. Quake employs a multi-level partitioning scheme that adjusts to updates and changing access patterns, guided by a cost model that predicts query latency based on partition sizes and access frequencies. Quake also dynamically sets query execution parameters to meet recall targets using a novel recall estimation model. Furthermore, Quake utilizes NUMA-aware intra-query parallelism for improved memory bandwidth utilization during search. To evaluate Quake, we prepare a Wikipedia vector search workload and develop a workload generator to create vector search workloads with configurable access patterns. Our evaluation shows that on dynamic workloads, Quake achieves query latency reductions of $1.5$–$38\times$ and update latency reductions of $4.5$–$126\times$ compared to state-of-the-art indexes such as SVS, DiskANN, HNSW, and SCANN.